%============================================================
%  appendix_c.tex  –  Implementation and Code
%============================================================
\chapter{Implementation Details and Source Code}
\label{app:code}

% ---------------------------------------------------------------
% WRITING TIPS – APPENDIX C
% ---------------------------------------------------------------
% Purpose: Provide enough implementation detail that the work is
%           fully reproducible. At UCPH, you do not need to include
%           every line, but key algorithmic functions should be shown.
%
% Suggested content:
%
%  C.1  Repository and Version Information
%       - Git repository URL (or note that code is available on request).
%       - Commit hash of the version used to produce thesis results.
%       - Python and PySCF versions (repeat from Computational Details).
%
%  C.2  Key Implementation Choices
%       - Explain any major departures from the "textbook" algorithm.
%       - E.g., use of interleaved α/β spin-orbital ordering vs.
%         blocked ordering.
%       - How the gradient and Hessian indices are mapped to flat arrays.
%       - Level-shifting strategy (code snippet + explanation).
%
%  C.3  Core Functions
%       - Include listings (using the listings package, style=python) of:
%         * MP2_energy() – the amplitude and energy calculation.
%         * orbital_optimization() – gradient, Hessian, Newton step.
%         * spatial_to_spin_eri() – the integral conversion.
%       - Annotate the code with comments referencing equation numbers
%         from the Theory chapter.
%
%  C.4  Profiling and Performance Notes
%       - Summary of cProfile results (from branch/profil/).
%       - Table: function | calls | cumtime | % of total runtime.
%       - Recommendations (even if not implemented) for future speedup.
%
% Tips:
%   - Use \lstset{style=python} and the \begin{lstlisting} environment.
%   - Keep code listings < 60 lines each; split into meaningful units.
%   - Comment the code in the listing if the thesis comments are better
%     than the source file comments.
%   - Do NOT put the entire ovos.py here – select the most important
%     400–500 lines.
%   - Acknowledge AI-assisted code writing per UCPH guidelines.
% ---------------------------------------------------------------

\section{Repository and Version Information}

% Example: point to your Git repo
% \noindent Code available at:
% \url{https://github.com/your-username/optimized_virtual_orbitals}\\
% Commit: \texttt{$\langle$hash$\rangle$}

\section{Spin-Orbital Integral Conversion}

\begin{lstlisting}[style=python, caption={Spatial-to-spin ERI conversion (\texttt{spatial\_2\_spin\_eri\_optimized}).}, label=lst:eri_conv]
# Convert spatial ERIs to spin-orbital basis (interleaved alpha/beta)
def spatial_2_spin_eri_optimized(self, eri_aaaa, eri_aabb, eri_bbbb):
    n_spin = self.tot_num_spin_orbs
    eri_spin = np.zeros((n_spin, n_spin, n_spin, n_spin))
    # alpha-alpha block (even indices)
    eri_spin[0::2, 0::2, 0::2, 0::2] = eri_aaaa
    # beta-beta block (odd indices)
    eri_spin[1::2, 1::2, 1::2, 1::2] = eri_bbbb
    # alpha-alpha / beta-beta cross block
    eri_spin[0::2, 0::2, 1::2, 1::2] = eri_aabb
    eri_spin[1::2, 1::2, 0::2, 0::2] = eri_aabb.transpose(2, 3, 0, 1)
    return eri_spin
\end{lstlisting}

\section{MP2 Energy and Amplitude Computation}

\section{Orbital Optimisation: Gradient and Hessian}

\section{Profiling Results}
